% frontmatter.tex





\thispagestyle{empty}
\begin{center}
	\centering
  \includegraphics{figures/uq_logo}

  \vspace{50pt}

  \textbf{\Large Trajectory Optimisation of a Partially-Reusable Rocket-Scramjet-Rocket Small Satellite Launch System}

  \vspace{25pt}

  {\large Sholto O. Forbes-Spyratos }
    
    \vspace{5pt}

  {\large B.Eng. (Mechanical and Aerospace) (Hons. I) \& B.Sc. (Physics)}


  \vspace{60pt}
                                                                                                                                                                                                                                                                                                                                                                                                                                                                                                                                                                                                                                                                                                                                                                                                                                                                                                                                                                                                                                                 
  \vfill

  {\large A thesis submitted for the degree of Doctor of Philosophy at\\ The University of Queensland in 2020}

  \vspace{20pt}

  School of Mechanical and Mining Engineering
  
  \textit{Centre for Hypersonics}

  \vfill
\end{center}

\cleardoublepage
\textcolor{black}{
\section*{Abstract}
}
\noindent
 The small satellite industry is expanding rapidly, driving a need for dedicated and cost-effective small satellite launchers.  
These small satellite launchers are currently designed to be single-use and low-cost, however; significant improvements in cost effectiveness and responsiveness may be achieved via the introduction of reusability. Unfortunately, the inclusion of reusability into small launch system designs is significantly more challenging than for larger launchers, and novel design methods may be necessary to achieve reliable and efficient re-use. One possibility for accomplishing reusability in a small launch system is the inclusion of airbreathing propulsion systems.

Airbreathing propulsion systems are extremely efficient, and do not require oxidiser to be carried on-board a launch system, bringing considerable operational and design advantages. 
However; airbreathing engines also add complexity to a launch system, with intricate subsystems and limited flight regimes. 
Multi-stage systems are likely to be the most applicable to small launch systems, but these designs compound on the complexity of airbreathing launch systems, and the efficiency trade-offs between the stages, particularly for three-stage systems, are not well understood. This work develops and analyses optimal payload-to-orbit trajectories for a three-stage airbreathing launcher, in order to understand the capabilities of the system and investigate the performance trade-offs between the stages during flight.



 A representative launch system is developed, consisting of two rocket stages and scramjet-powered accelerator stage.
This launch system is based on the SPARTAN scramjet-powered accelerator, being developed by the University of Queensland and Hypersonix, with a first stage based on the Falcon-1e, and a bespoke third stage. These stages are designed to the detail required for preliminary analysis, along with aerodynamic, mass and propulsive models.
 A package is developed to calculate the maximum payload-to-orbit trajectory of airbreathing launch systems, designated LODESTAR. LODESTAR utilises GPOPS-2, a pseudospectral method optimal control software, configured to calculate maximum payload-to-orbit trajectory profiles. Along with the configuration of GPOPS-2, LODESTAR provides a three degree-of-freedom simulation of each stage, and tools to verify and examine the optimised solutions produced by GPOPS-2.


 
Launch trajectories are initially simulated assuming that the scramjet accelerator lands at some position downrange. A launch trajectory is simulated in which the scramjet accelerator is constrained to flight at maximum dynamic pressure for reference and verification, as flight at this limit during airbreathing engine operation is often assumed in previous studies to produce a maximum efficiency launch trajectory. This trajectory achieves a payload-to-orbit of \PayloadToOrbitConstqNoReturn kg, launching to sun synchronous orbit.  
The optimised maximum efficiency trajectory of the launch system is then calculated, and is found to differ significantly from the trajectory in which the scramjet accelerator is constrained to constant dynamic pressure. 
 The scramjet accelerator is found to deviate from its maximum dynamic pressure at both stage separation points, and for a segment in the middle of its trajectory.
 The higher separation points result in the efficiency of the scramjet accelerator reducing, but increase the efficiency of the rocket stages, improving the overall efficiency of the system. 
  Additionally, an altitude raising manoeuvre is performed in a region where the specific impulse of the scramjet engines is relatively homogeneous with varied flight conditions, resulting in a further small performance increase. Overall, flying an optimal maximum efficiency trajectory increases the payload-to-orbit of the system launching to sun synchronous orbit to \PayloadToOrbitStandardNoReturn kg, an increase of 59.1\% compared to a trajectory in which the scramjet accelerator flies at a constant dynamic pressure.  
 
 
 Next, the fly-back of the scramjet accelerator is included within the trajectory optimisation, and a maximum payload-to-orbit flight profile is calculated. 
 It is found that the scramjet accelerator is able to achieve fly-back, but that it must ignite its scramjet engines during its return flight, causing the fly-back to become an important consideration in the optimal trajectory design. When the fly-back is included, the first stage pitches easterly, rather than northerly. The scramjet accelerator banks heavily throughout its acceleration to manoeuvre to polar inclination, decreasing the performance of the scramjet accelerator, but also reducing the amount of fuel used during fly-back, for a net performance gain. 
The fly-back is found to exhibit multiple `skipping' manoeuvres, which serve to increase the range of the scramjet accelerator, minimising the fuel necessary during the return flight. In addition, the scramjet engines are powered on at the troughs of the first three skips, corresponding to the points of highest possible specific impulse. This optimal trajectory for the launch system is able to deliver \PayloadToOrbitStandard kg of payload to sun synchronous orbit while successfully returning the scramjet accelerator to its initial launch site.  

Finally, a study is conducted to quantify the sensitivity of the launch system to variations in key design parameters. 
The sensitivities of coupled design parameters are compared, to quantify their relative impacts on the performance of the launch system. The magnitudes of these relative impacts are assessed, and the behaviour of the maximum payload-to-orbit trajectory is investigated as the performance characteristics of the launch system are modified, to provide insights into both the optimised trajectory and the launch vehicle performance. 

 

\clearpage
\section*{Declaration by author}
This thesis is composed of my original work, and contains no material previously published or written by another person except where due reference has been made in the text. I have clearly stated the contribution by others to jointly-authored works that I have included in my thesis.

I have clearly stated the contribution of others to my thesis as a whole, including statistical assistance, survey design, data analysis, significant technical procedures, professional editorial advice, financial support and any other original research work used or reported in my thesis. The content of my thesis is the result of work I have carried out since the commencement of my higher degree by research candidature and does not include a substantial part of work that has been submitted to qualify for the award of any other degree or diploma in any university or other tertiary institution. I have clearly stated which parts of my thesis, if any, have been submitted to qualify for another award.

I acknowledge that an electronic copy of my thesis must be lodged with the University Library and, subject to the policy and procedures of The University of Queensland, the thesis be made available for research and study in accordance with the Copyright Act 1968 unless a period of embargo has been approved by the Dean of the Graduate School. 

I acknowledge that copyright of all material contained in my thesis resides with the copyright holder(s) of that material. Where appropriate I have obtained copyright permission from the copyright holder to reproduce material in this thesis and have sought permission from co-authors for any jointly authored works included in the thesis.


\clearpage
\section*{Publications Included in this Thesis}

This thesis comprises partly of publications, as allowed by University of Queensland Policy PPL 4.60.07. The papers that have been included have all been published in peer reviewed journals at the time of submission. 

\vspace{\baselineskip}
\noindent
\AtNextCite{\defcounter{maxnames}{99}}\fullcite{ForbesSpyratos2018}

\begin{center}
  \begin{tabular}{ll}
    \toprule
    Contributor   & Contribution \\
    \midrule
    Sholto O. Forbes-Spyratos             
                                  & Conception and Design (85\%)\\
                                  & Performed simulations (100\%)\\
                                  & Analysis of results (90\%)\\
                                  & Wrote and edited paper (85\%)\\
    \midrule
    Ingo H. Jahn            
                                  & Conception and Design (5\%)\\
                                  & Analysis of results (5\%)\\
                                  & Wrote and edited paper (7.5\%)\\
                                  
        \midrule
        Michael P. Kearney            
								      & Conception and Design (5\%)\\
								      & Wrote and edited paper (7.5\%)\\                       

    \midrule
    Michael K. Smart              
                                  & Conception and Design (5\%)\\
                                  & Analysis of results (5\%)\\
                                  & Wrote and edited paper (5\%)\\
    \bottomrule
  \end{tabular}
\end{center}

\section*{Other Publications During Candidature}

\subsection*{Conference papers}

\noindent\AtNextCite{\defcounter{maxnames}{99}}\fullcite{ForbesSpyratos2017}\\

\noindent\AtNextCite{\defcounter{maxnames}{99}}\fullcite{Forbes2018a}\\

\noindent\AtNextCite{\defcounter{maxnames}{99}}\fullcite{chai2017}\\

\section*{Contributions by Others to the Thesis}

The initial model of the Baseline SPARTAN was provided for this work by Dr. Dawid Preller and Mr. Joseph Chai, including mass properties, dimensions, and CAD models. The \textsf{CRESTM10} scramjet engine database, consisting of tabulated performance data over a range of inlet conditions, and the \textsf{cone\_shoot} program were provided for this study by Prof. Michael Smart. The viscous correction incorporated into the aerodynamic calculations and the coefficient of friction calculations incorporated into the thermal analysis were performed by Mr. Alexander Ward in collaboration with the Author, and provided for this study in the form of tabulated databases. The initial \textsf{SurfaceTemp} code that was modified for the thermal analysis in this work was provided by Dr. Ingo Jahn. 



\section*{Submitted Manuscripts Included in this Thesis}
No manuscripts included.

\section*{Research Involving Human or Animal Subjects}
No animal or human subjects were involved in this research.

\section*{Statement of Parts of the Thesis Submitted to Qualify for the Award of Another Degree}

No works submitted towards another degree have been included in this thesis.

\clearpage

\section*{Financial Support}
This research was supported by an Australian Government Research Training Program Scholarship and a University of Queensland Candidate Development Award.
\clearpage


\section*{Acknowledgements} 
  For Kaitlin, with love to my parents Alexandra and Tony, and utmost gratitude to my family and friends, as well as to my advisors; Ingo Jahn, Michael Kearney, and Michael Smart. 

      \vspace*{\fill}


\clearpage
\subsection*{Keywords}
  airbreathing propulsion, scramjets, hypersonics, access-to-space, small satellite launchers, airbreathing launch systems, trajectory optimisation

\subsection*{Australian and New Zealand Standard Research Classification (ANZSRC)}

ANZSRC code: 090107 Hypersonic Propulsion and Hypersonic Aerodynamics, 30\% \newline
ANZSRC code: 090108, Satellite, Space Vehicle and Missile Design and Testing, 40\% \newline
ANZSRC code: 010303 Optimisation, 30\% \newline

\subsection*{Fields of Research (FoR) Classification}

 FoR code: 0901, Aerospace Engineering, 100\%

\tableofcontents

\listoffigures
\addcontentsline{toc}{chapter}{List of Figures}
\listoftables
\addcontentsline{toc}{chapter}{List of Tables}

\newpage
\section*{Nomenclature}

\addcontentsline{toc}{chapter}{Nomenclature}
\markboth{NOMENCLATURE}{NOMENCLATURE}


\noindent
\begin{tabular}{p{0.8cm}p{0.8cm}p{5.6cm}}
	$a$ & $=$ & Acceleration (m/s$^2$)\\
\end{tabular} 
\begin{tabular}{p{0.8cm}p{0.8cm}p{5.6cm}}
	$C_L,C_D$ & $=$ & Aerodynamic Coefficients\\
\end{tabular} 
\begin{tabular}{p{0.8cm}p{0.8cm}p{5.6cm}}
	$\alpha$ & $=$ & Angle of Attack (rad)\\
\end{tabular} 
\begin{tabular}{p{0.8cm}p{0.8cm}p{5.6cm}}
	$\omega$ & $=$ & Angular Velocity (rad/s)\\
\end{tabular} 
\begin{tabular}{p{0.8cm}p{0.8cm}p{5.6cm}}
	$b$ & $=$ & Bound Constraint Values\\
\end{tabular}
\begin{tabular}{p{0.8cm}p{0.8cm}p{5.6cm}}
	$w_{cap}$ & $=$ & Capture Width\\
\end{tabular} 
\begin{tabular}{p{0.8cm}p{0.8cm}p{5.6cm}}
	$\textbf{u}$& $=$ & Control Variables\\
\end{tabular} 
\begin{tabular}{p{0.8cm}p{0.8cm}p{5.6cm}}
	$\lambda$ & $=$ & Costate Variablraint 
\end{tabular}
\begin{tabular}{p{0.8cm}p{0.8cm}p{5.6cm}}
	$\rho$ & $=$ & Density (kg/m$^2$)\\
\end{tabular} 
\begin{tabular}{p{0.8cm}p{0.8cm}p{5.6cm}}
	$D$ & $=$ & Drag (N)\\
\end{tabular} 
\begin{tabular}{p{0.8cm}p{0.8cm}p{5.6cm}}
	$\eta$ & $=$ & Efficiency\\
\end{tabular} 
\begin{tabular}{p{0.8cm}p{0.8cm}p{5.6cm}}
	$q$ & $=$ & Dynamic Pressure (Pas)\\  
\end{tabular}     
\begin{tabular}{p{0.8cm}p{0.8cm}p{5.6cm}}
	$\psi$ & $=$ & Equality Constraint Function\\
\end{tabular}
\begin{tabular}{p{0.8cm}p{0.8cm}p{5.6cm}}
	$\gamma$ & $=$ & Flight Path Angle (rad)\\
\end{tabular} 
\begin{tabular}{p{0.8cm}p{0.8cm}p{5.6cm}}
	$F$ & $=$ & Force (N)\\
\end{tabular} 
\begin{tabular}{p{0.8cm}p{0.8cm}p{5.6cm}}
	$R$ & $=$ & Gas Constant (8.314 $J K^{-1} mol^{-1}$)\\
\end{tabular} 
\begin{tabular}{p{0.8cm}p{0.8cm}p{5.6cm}}
	$g$ & $=$ & Gravitational Acceleration at Earth's Surface (m/s$^2$)\\
\end{tabular}
\begin{tabular}{p{0.8cm}p{0.8cm}p{5.6cm}}
	$H$ & $=$ & Hamiltonian\\
\end{tabular}
\begin{tabular}{p{0.8cm}p{0.8cm}p{5.6cm}}
	$\zeta$ & $=$ & Heading Angle (rad)\\
\end{tabular}
\begin{tabular}{p{0.8cm}p{0.8cm}p{5.6cm}}
	$\lambda$ & $=$ & Inequality Constraint Function\\
\end{tabular}
\begin{tabular}{p{0.8cm}p{0.8cm}p{5.6cm}}
	$KE$ & $=$ & Kinetic Energy (J)\\
\end{tabular} 
\begin{tabular}{p{0.8cm}p{0.8cm}p{5.6cm}}
	$\phi$ & $=$ & Latitude (rad) \& Equivalence Ratio \& Heat Flux (W/m$^2$)\\
\end{tabular}
\begin{tabular}{p{0.8cm}p{0.8cm}p{5.6cm}}
	$L$ & $=$ & Lift (N) \& Integrated Cost Function\\
\end{tabular} 
\begin{tabular}{p{0.8cm}p{0.8cm}p{5.6cm}}
	$\xi$ & $=$ & Longitude (rad)\\
\end{tabular} 
\begin{tabular}{p{0.8cm}p{0.8cm}p{5.6cm}}
	$M$ & $=$ & Mach Number \& Endpoint Cost Function\\
\end{tabular}
\begin{tabular}{p{0.8cm}p{0.8cm}p{5.6cm}}
	$m$ & $=$ & Mass (kg)\\
\end{tabular} 
\begin{tabular}{p{0.8cm}p{0.8cm}p{5.6cm}}
	$PE$ & $=$ & Potential Energy (J)\\
\end{tabular} 
\begin{tabular}{p{0.8cm}p{0.8cm}p{5.6cm}}
	$P$ & $=$ & Pressure (Pa)  \\
\end{tabular}
\begin{tabular}{p{0.8cm}p{0.8cm}p{5.6cm}}
	$r$ & $=$ & Radius from Earth Centre (m)\\
\end{tabular} 
\begin{tabular}{p{0.8cm}p{0.8cm}p{5.6cm}}
	$A$ & $=$ & Reference Area (m$^2$)\\
\end{tabular} 
\begin{tabular}{p{0.8cm}p{0.8cm}p{5.6cm}}
	$I_{sp}$ & $=$ & Specific Impulse (s)\\ 
\end{tabular} 
\begin{tabular}{p{0.8cm}p{0.8cm}p{5.6cm}}
	$v$ & $=$ & Speed (m/s)\\
\end{tabular} 
\begin{tabular}{p{0.8cm}p{0.8cm}p{5.6cm}}
	$\textbf{x}$& $=$ & State Variables\\
\end{tabular} 
\begin{tabular}{p{0.8cm}p{0.8cm}p{5.6cm}}
	$T$ & $=$ & Thrust (N) \& Temperature (K)\\
\end{tabular} 
\begin{tabular}{p{0.8cm}p{0.8cm}p{5.6cm}}
	$t$ & $=$ & Time (s)\\
\end{tabular} 
\begin{tabular}{p{0.8cm}p{0.8cm}p{5.6cm}}
	$J$ & $=$ & Total Cost Functional\\
\end{tabular}
\newline  	
\begin{tabular}{p{5.2cm}p{1cm}p{10cm}}
	\textit{Subscripts} & & \\
\end{tabular} 
\begin{tabular}{p{0.8cm}p{0.8cm}p{5.6cm}}
	$1$ & $=$ & $1^{st}$ Stage Rocket\\
\end{tabular} 
\begin{tabular}{p{0.8cm}p{0.8cm}p{5.6cm}}
	$2$ & $=$ & $2^{nd}$ Stage Scramjet Vehicle\\
\end{tabular} 
\begin{tabular}{p{0.8cm}p{0.8cm}p{5.6cm}}
	$3$ & $=$ & $3^{rd}$ Stage Rocket\\
\end{tabular} 
\begin{tabular}{p{0.8cm}p{0.8cm}p{5.6cm}}
	$\rightarrow$ & = & Stage Transition\\
\end{tabular} 
\begin{tabular}{p{0.9cm}p{0.7cm}p{5.6cm}}
	\rom{1}-\rom{7} & $=$ & Trajectory Phases\\
\end{tabular} 
\begin{tabular}{p{0.8cm}p{0.8cm}p{5.6cm}}
	AoA & $=$ & Angle of Attack\\
\end{tabular} 
\begin{tabular}{p{0.8cm}p{0.8cm}p{5.6cm}}
	$atm$ & $=$ & Atmospheric\\
\end{tabular} 
\begin{tabular}{p{0.8cm}p{0.8cm}p{5.6cm}}
	$E$ & $=$ & Earth\\
\end{tabular} 
\begin{tabular}{p{0.8cm}p{0.8cm}p{5.6cm}}
	e & $=$ & Exit\\
\end{tabular} 
\begin{tabular}{p{0.8cm}p{0.8cm}p{5.6cm}}
	$f$ & $=$ & Final\\
\end{tabular}
\begin{tabular}{p{0.8cm}p{0.8cm}p{5.6cm}}
	$0$ & $=$ & Initial\\
\end{tabular} 
 \begin{tabular}{p{0.8cm}p{0.8cm}p{5.6cm}}
 	LH2 & $=$ & Liquid Hydrogen\\
 \end{tabular} 
\begin{tabular}{p{0.8cm}p{0.8cm}p{5.6cm}}
	LOX & $=$ & Liquid Oxygen\\
\end{tabular} 
\begin{tabular}{p{0.8cm}p{0.8cm}p{5.6cm}}
	$LEO$ & $=$ & Low Earth Orbit\\
\end{tabular}
\begin{tabular}{p{0.8cm}p{0.8cm}p{5.6cm}}
	$a,b$ & $=$ & Nominal Stage Numbers\\
\end{tabular} 
\begin{tabular}{p{0.8cm}p{0.8cm}p{5.6cm}}
	$p$ & $=$ & Propulsive\\
\end{tabular} 
\begin{tabular}{p{0.8cm}p{0.8cm}p{5.6cm}}
	$SL$ & $=$ & Sea Level\\
\end{tabular} 
\begin{tabular}{p{0.8cm}p{0.8cm}p{5.6cm}}
	$vac$ & $=$ & Vacuum\\
\end{tabular} 
\begin{tabular}{p{0.8cm}p{0.8cm}p{5.6cm}}
	$r$ & $=$ & Radial\\
\end{tabular} 
\begin{tabular}{p{0.8cm}p{0.8cm}p{5.6cm}}
	$t$ & $=$ & Tangential\\
\end{tabular} 

